% Auto-generated from raw.txt
% Usage:
%   \vocabentry{word}{/ipa/ (pos)}{definition}{example}{note}

\vocabentry{Heritage}
{/ˈherɪtɪdʒ/ (n)}
{Property that is or may be inherited; an inheritance.}
{The city is working hard to preserve its rich architectural heritage from the colonial era.}
{Đây là một danh từ dùng để chỉ di sản, bao gồm các giá trị văn hóa, truyền thống hoặc công trình lịch sử để lại.}

\vocabentry{Artifact}
{/ˈɑːrtɪfækt/ (n)}
{An object made by a human being, typically one of cultural or historical interest.}
{Ancient artifacts discovered in the tomb provide clues about the daily life of the Pharaohs.}
{Đây là một danh từ dùng để chỉ hiện vật, những đồ vật cổ do con người tạo ra.}

\vocabentry{Archaeology}
{/ˌɑːrkiˈɑːlədʒi/ (n)}
{The study of human history and prehistory through the excavation of sites.}
{Archaeology allows us to reconstruct past civilizations by examining physical remains.}
{Đây là một danh từ dùng để chỉ khảo cổ học.}

\vocabentry{Chronology}
{/krəˈnɑːlədʒi/ (n)}
{The arrangement of events or dates in the order of their occurrence.}
{The historian established a clear chronology of the events leading up to the revolution.}
{Đây là một danh từ dùng để chỉ niên đại học hoặc sự sắp xếp theo thứ tự thời gian.}

\vocabentry{Civilization}
{/ˌsɪvələˈzeɪʃn/ (n)}
{The stage of human social and cultural development and organization that is considered most advanced.}
{Ancient Greek civilization had a profound influence on Western art and philosophy.}
{Đây là một danh từ dùng để chỉ nền văn minh.}

\vocabentry{Colonialism}
{/kəˈloʊniəlɪzəm/ (n)}
{The policy or practice of acquiring full or partial political control over another country.}
{The legacy of colonialism is still evident in the political structures of many African nations.}
{Đây là một danh từ dùng để chỉ chủ nghĩa thực dân.}

\vocabentry{Era}
{/ˈɪrə/ (n)}
{A long and distinct period of history with a particular feature or characteristic.}
{The Victorian era was characterized by rapid industrialization and social change.}
{Đây là một danh từ dùng để chỉ một kỷ nguyên hoặc một thời đại lịch sử.}

\vocabentry{Dynasty}
{/ˈdaɪnəsti/ (n)}
{A line of hereditary rulers of a country.}
{The Ming Dynasty is famous for its exquisite porcelain and the building of the Great Wall.}
{Đây là một danh từ dùng để chỉ triều đại, một dòng họ cầm quyền liên tục.}

\vocabentry{Ancestor}
{/ˈænsestər/ (n)}
{A person, typically one more remote than a grandparent, from whom one is descended.}
{Many people travel to their ancestral villages to learn more about their family history.}
{Đây là một danh từ dùng để chỉ tổ tiên.}

\vocabentry{Preservation}
{/ˌprezərˈveɪʃn/ (n)}
{The action of preserving something; keeping it in its original state.}
{The preservation of ancient manuscripts is vital for our understanding of medieval literature.}
{Đây là một danh từ dùng để chỉ sự bảo tồn nguyên trạng các giá trị lịch sử.}

\vocabentry{Authenticity}
{/ˌɔːθenˈtɪsəti/ (n)}
{The quality of being authentic or genuine.}
{Experts were called in to verify the authenticity of the painting attributed to Rembrandt.}
{Đây là một danh từ dùng để chỉ tính xác thực hoặc nguyên bản của một di tích/hiện vật.}

\vocabentry{Folklore}
{/ˈfoʊklɔːr/ (n)}
{The traditional beliefs, customs, and stories of a community, passed through the generations.}
{The author drew heavy inspiration from local folklore when writing his historical novels.}
{Đây là một danh từ dùng để chỉ văn hóa dân gian hoặc truyền thuyết truyền miệng.}

\vocabentry{Imperialism}
{/ɪmˈpɪriəlɪzəm/ (n)}
{A policy of extending a country's power and influence through diplomacy or military force.}
{19th-century imperialism led to the partition of Africa among European powers.}
{Đây là một danh từ dùng để chỉ chủ nghĩa đế quốc.}

\vocabentry{Legacy}
{/ˈleɡəsi/ (n)}
{An amount of money or property left to someone in a will; or the long-lasting impact of events.}
{The philosopher's legacy continues to influence political thought to this day.}
{Đây là một danh từ dùng để chỉ di sản hoặc hệ quả kéo dài của quá khứ.}

\vocabentry{Monument}
{/ˈmɑːnjumənt/ (n)}
{A statue, building, or other structure erected to commemorate a notable person or event.}
{The monument was built to honor those who lost their lives in the Great War.}
{Đây là một danh từ dùng để chỉ đài tưởng niệm hoặc công trình vĩ đại.}

\vocabentry{Prehistoric}
{/ˌpriːhɪˈstɔːrɪk/ (adj)}
{Relating to or denoting the period before written records.}
{Prehistoric cave paintings show that art has been part of human life for millennia.}
{Đây là một tính từ dùng để chỉ thời kỳ tiền sử.}

\vocabentry{Revival}
{/rɪˈvaɪvl/ (n)}
{An improvement in the condition or strength of something; an instance of something becoming popular again.}
{There has been a recent revival of interest in traditional folk music.}
{Đây là một danh từ dùng để chỉ sự phục hưng hoặc sự sống lại của một truyền thống.}

\vocabentry{Sovereignty}
{/ˈsɑːvrənti/ (n)}
{Supreme power or authority; the authority of a state to govern itself.}
{The treaty was signed to protect the sovereignty and territorial integrity of the nation.}
{Đây là một danh từ dùng để chỉ chủ quyền quốc gia.}

\vocabentry{Tradition}
{/trəˈdɪʃn/ (n)}
{The transmission of customs or beliefs from generation to generation.}
{It is a long-standing tradition in our family to gather for a big meal on New Year's Eve.}
{Đây là một danh từ dùng để chỉ truyền thống.}

\vocabentry{Excavation}
{/ˌekskəˈveɪʃn/ (n)}
{The action of excavating something, especially an archaeological site.}
{The excavation of the ancient city revealed a complex system of tunnels and chambers.}
{Đây là một danh từ dùng để chỉ việc khai quật khảo cổ.}

\vocabentry{Antique}
{/ænˈtiːk/ (n/adj)}
{A collectible object such as a piece of furniture or work of art that has a high value because of its age.}
{He owns a valuable collection of antique clocks from the 18th century.}
{Đây là một danh từ/tính từ chỉ đồ cổ hoặc thuộc về thời xưa.}

\vocabentry{Archive}
{/ˈɑːrkaɪv/ (n/v)}
{A collection of historical documents or records providing information about a place, institution, or group of people.}
{Historians spent months researching in the national archives to find the truth about the event.}
{Đây là một danh từ/động từ chỉ kho lưu trữ tư liệu.}

\vocabentry{Feudalism}
{/ˈfjuːdlɪzəm/ (n)}
{The dominant social system in medieval Europe.}
{Feudalism was based on the relationship between lords and vassals.}
{Đây là một danh từ dùng để chỉ chế độ phong kiến.}

\vocabentry{Inscription}
{/ɪnˈskrɪpʃn/ (n)}
{Words inscribed, as on a monument or in a book.}
{The stone tablet bore an ancient inscription that was difficult to translate.}
{Đây là một danh từ dùng để chỉ văn bia hoặc câu chữ khắc trên hiện vật.}

\vocabentry{Medieval}
{/ˌmediˈiːvl/ (adj)}
{Relating to the Middle Ages.}
{Medieval knights were expected to follow a strict code of conduct called chivalry.}
{Đây là một tính từ dùng để chỉ thời Trung cổ, khoảng thế kỷ 5-15}

\vocabentry{Oral history}
{/ˈɔːrəl ˈhɪstri/ (n)}
{The collection and study of historical information using sound recordings of interviews with people having personal knowledge of past events.}
{Oral history provides a unique perspective on the lives of ordinary people in the past.}
{Đây là một cụm danh từ dùng để chỉ sử học truyền miệng.}

\vocabentry{Renaissance}
{/ˌrenəˈsɑːns/ (n)}
{The revival of art and literature under the influence of classical models in the 14th–16th centuries.}
{The Renaissance marked the transition from the Middle Ages to modernity.}
{Đây là một danh từ dùng để chỉ thời kỳ Phục hưng.}

\vocabentry{Treaty}
{/ˈtriːti/ (n)}
{A formally concluded and ratified agreement between countries.}
{The peace treaty brought an end to decades of conflict between the two nations.}
{Đây là một danh từ dùng để chỉ hiệp ước hoặc hòa ước.}

\vocabentry{Uprising}
{/ˈʌpraɪzɪŋ/ (n)}
{An act of resistance or rebellion; a revolt.}
{The popular uprising eventually led to the downfall of the monarchy.}
{Đây là một danh từ dùng để chỉ cuộc khởi nghĩa hoặc nổi dậy.}

\vocabentry{Relic}
{/ˈrelɪk/ (n)}
{An object surviving from an earlier time, especially one of historical or religious interest.}
{The museum houses a holy relic believed to have belonged to a famous saint.}
{Đây là một danh từ dùng để chỉ di tích, di vật hoặc thánh tích.}

\vocabentry{Nationalism}
{/ˈnæʃnəlɪzəm/ (n)}
{Identification with one's own nation and support for its interests.}
{Nationalism can be a powerful force for national unity during times of crisis.}
{Đây là một danh từ dùng để chỉ chủ nghĩa dân tộc, tư tưởng yêu nước và ủng hộ quyền lợi quốc gia}

\vocabentry{Propaganda}
{/ˌprɑːpəˈɡændə/ (n)}
{Information, especially of a biased or misleading nature, used to promote a political cause or point of view.}
{The regime used propaganda to maintain control and suppress political dissent.}
{Đây là một danh từ dùng để chỉ sự tuyên truyền.}

\vocabentry{Revolution}
{/ˌrevəˈluːʃn/ (n)}
{A forcible overthrow of a government or social order, in favor of a new system.}
{The French Revolution had a profound impact on the course of modern history.}
{Đây là một danh từ dùng để chỉ cuộc cách mạng.}

\vocabentry{Democracy}
{/dɪˈmɑːkrəsi/ (n)}
{A system of government by the whole population or all the eligible members of a state.}
{Ancient Athens is often cited as the birthplace of Western democracy.}
{Đây là một danh từ dùng để chỉ chế độ dân chủ.}

\vocabentry{Monarchy}
{/ˈmɑːnərki/ (n)}
{A form of government with a monarch at the head.}
{The French monarchy was abolished during the revolution of 1789.}
{Đây là một danh từ dùng để chỉ chế độ quân chủ, nơi quyền lực cao nhất thuộc về một vị vua hoặc nữ hoàng}

\vocabentry{Primary source}
{/ˈpraɪmeri sɔːrs/ (n)}
{An artifact, document, diary, manuscript, autobiography, recording, or any other source of information that was created at the time under study.}
{Letters written by soldiers during the war are valuable primary sources for historians.}
{Đây là một cụm danh từ dùng để chỉ nguồn tư liệu sơ cấp.}

\vocabentry{Secondary source}
{/ˈsekənderi sɔːrs/ (n)}
{A document or recording that relates or discusses information originally presented elsewhere.}
{A modern history textbook is a secondary source that interprets past events.}
{Đây là một cụm danh từ dùng để chỉ nguồn tư liệu thứ cấp.}

\vocabentry{Indigenous}
{/ɪnˈdɪdʒənəs/ (adj)}
{Originating or occurring naturally in a particular place; native.}
{The indigenous people of the island have lived there for thousands of years.}
{Đây là một tính từ dùng để chỉ bản địa, thuộc về người dân hoặc văn hóa có nguồn gốc từ một vùng đất cụ thể}

\vocabentry{Patrimony}
{/ˈpætrɪmoʊni/ (n)}
{Property inherited from one's father or male ancestor; or the heritage of a nation.}
{The historical monuments are part of the nation's cultural patrimony and must be protected.}
{Đây là một danh từ trang trọng dùng để chỉ tài sản gia truyền hoặc di sản quốc gia.}

\vocabentry{Restoration}
{/ˌrestəˈreɪʃn/ (n)}
{The action of returning something to a former owner, place, or condition.}
{The restoration of the old opera house took several years and cost millions of dollars.}
{Đây là một danh từ dùng để chỉ sự trùng tu hoặc phục hồi di tích.}

\vocabentry{Ancestry}
{/ˈænsestri/ (n)}
{One's family or ethnic descent.}
{He is proud of his Scottish ancestry and often wears a kilt on special occasions.}
{Đây là một danh từ dùng để chỉ dòng dõi hoặc nguồn gốc tổ tiên.}

\vocabentry{Antiquity}
{/ænˈtɪkwəti/ (n)}
{The ancient past, especially the period before the Middle Ages.}
{Many of the world's greatest philosophers lived in classical antiquity.}
{Đây là một danh từ dùng để chỉ thời cổ đại hoặc những vật cổ xưa.}

\vocabentry{Bicentennial}
{/ˌbaɪsenˈteniəl/ (n/adj)}
{The two-hundredth anniversary of a significant event.}
{The city held a massive parade to celebrate its bicentennial.}
{Đây là một danh từ/tính từ dùng để chỉ kỷ niệm 200 năm.}

\vocabentry{Centuries}
{/ˈsentʃəriz/ (n)}
{Periods of one hundred years.}
{The building has stood for several centuries and has seen many historical events.}
{Đây là một danh từ dùng để chỉ các thế kỷ.}

\vocabentry{Conflict}
{/ˈkɑːnflɪkt/ (n)}
{A serious disagreement or argument; a prolonged armed struggle.}
{Historians analyze the causes and consequences of major international conflicts.}
{Đây là một danh từ dùng để chỉ cuộc xung đột hoặc chiến tranh.}

\vocabentry{Conquest}
{/ˈkɑːŋkwest/ (n)}
{The subjugation and assumption of control of a place or people by use of military force.}
{The Norman Conquest of England in 1066 led to major changes in the country's laws and language.}
{Đây là một danh từ dùng để chỉ sự xâm chiếm hoặc cuộc chinh phạt.}

\vocabentry{Custodianship}
{/kʌˈstoʊdiənʃɪp/ (n)}
{The office or duty of a custodian; guardianship.}
{The local community was granted custodianship of the historic forest.}
{Đây là một danh từ dùng để chỉ quyền trông nom hoặc bảo hộ di sản.}

\vocabentry{Decade}
{/ˈdekeɪd/ (n)}
{A period of ten years.}
{The 1960s was a decade marked by significant social and cultural upheaval.}
{Đây là một danh từ dùng để chỉ thập kỷ.}

\vocabentry{Defeat}
{/dɪˈfiːt/ (n/v)}
{An instance of defeating or being defeated; to win a victory over.}
{The defeat of the Spanish Armada was a major turning point in English history.}
{Đây là một danh từ/động từ dùng để chỉ sự thất bại hoặc đánh bại đối thủ.}

\vocabentry{Descendant}
{/dɪˈsendənt/ (n)}
{A person, plant, or animal that is descended from a particular ancestor.}
{He claims to be a direct descendant of a famous explorer.}
{Đây là một danh từ dùng để chỉ hậu duệ hoặc con cháu đời sau.}

\vocabentry{Diplomacy}
{/dɪˈploʊməsi/ (n)}
{The profession, activity, or skill of managing international relations.}
{Diplomacy is often more effective than military force in resolving international disputes.}
{Đây là một danh từ dùng để chỉ ngành ngoại giao.}

\vocabentry{Discovery}
{/dɪˈskʌvəri/ (n)}
{The action or process of discovering or being discovered.}
{The discovery of the tomb of Tutankhamun in 1922 was a worldwide sensation.}
{Đây là một danh từ dùng để chỉ sự khám phá ra một vùng đất hoặc kiến thức mới.}

\vocabentry{Emperor}
{/ˈempərər/ (n)}
{A sovereign ruler of great power and rank, especially one ruling an empire.}
{Julius Caesar was one of the most famous emperors of the Roman Empire.}
{Đây là một danh từ dùng để chỉ hoàng đế.}

\vocabentry{Empire}
{/ˈempaɪər/ (n)}
{An extensive group of states or countries under a single supreme authority.}
{The British Empire was the largest in history, covering about a quarter of the world's land area.}
{Đây là một danh từ dùng để chỉ đế chế.}

\vocabentry{Enlightenment}
{/ɪnˈlaɪtnmənt/ (n)}
{A European intellectual movement of the late 17th and 18th centuries emphasizing reason and individualism.}
{The Enlightenment led to the development of modern science and democratic principles.}
{Đây là một danh từ dùng để chỉ thời kỳ Khai sáng.}

\vocabentry{Epoch}
{/ˈepək/ (n)}
{A period of time in history or a person's life, typically one marked by notable events or particular characteristics.}
{The fall of the Berlin Wall marked the end of a long epoch of European history.}
{Đây là một danh từ dùng để chỉ kỷ nguyên hoặc thời kỳ quan trọng.}

\vocabentry{Ethics}
{/ˈeθɪks/ (n)}
{Moral principles that govern a person's behavior or the conducting of an activity.}
{Historians must follow strict professional ethics to ensure the accuracy and objectivity of their work.}
{Đây là một danh từ dùng để chỉ đạo đức học, thường liên quan đến cách viết sử.}

\vocabentry{Ethnic}
{/ˈeθnɪk/ (adj)}
{Relating to a population group with a common national or cultural tradition.}
{The country is a melting pot of different ethnic groups, each with its own unique traditions.}
{Đây là một tính từ dùng để chỉ tính dân tộc hoặc sắc tộc.}

\vocabentry{Expansion}
{/ɪkˈspænʃn/ (n)}
{The action of becoming larger or more extensive.}
{The westward expansion of the United States in the 19th century led to many conflicts.}
{Đây là một danh từ dùng để chỉ sự mở rộng bờ cõi hoặc tầm ảnh hưởng.}

\vocabentry{Fable}
{/ˈfeɪbl/ (n)}
{A short story, typically with animals as characters, conveying a moral.}
{Aesop's fables are still widely read and used to teach moral lessons to children today.}
{Đây là một danh từ dùng để chỉ truyện ngụ ngôn mang tính giáo dục lịch sử.}

\vocabentry{Faction}
{/ˈfækʃn/ (n)}
{A small organized dissenting group within a larger one, especially in politics.}
{Rival factions within the government were struggling for control of the military.}
{Đây là một danh từ dùng để chỉ phe phái trong nội bộ chính quyền.}

\vocabentry{Folktale}
{/ˈfoʊkteɪl/ (n)}
{A story originating in popular culture, typically passed on by word of mouth.}
{Many folktales share similar themes and characters across different cultures.}
{Đây là một danh từ dùng để chỉ truyện dân gian.}

\vocabentry{Grave}
{/ɡreɪv/ (n)}
{A place of burial for a dead body, typically a hole dug in the ground and marked by a stone or mound.}
{Archaeologists discovered several ancient graves containing valuable gold jewelry.}
{Đây là một danh từ dùng để chỉ ngôi mộ, nguồn khảo cổ quan trọng.}

\vocabentry{Greatness}
{/ˈɡreɪtnəs/ (n)}
{The quality of being great; distinguished.}
{The historian wrote about the greatness and eventual fall of the Roman Empire.}
{Đây là một danh từ dùng để chỉ sự vĩ đại của một nhân vật hoặc đế chế.}

\vocabentry{Hearth}
{/hɜːrθ/ (n)}
{The floor of a fireplace; symbols of home and family.}
{In ancient times, the hearth was the center of social and family life.}
{Đây là một danh từ dùng để chỉ bếp lửa, biểu tượng cho gia đình và truyền thống trong sử học.}

\vocabentry{Hegemony}
{/hɪˈdʒeməni/ (n)}
{Leadership or dominance, especially by one country or social group over others.}
{The country's economic and military hegemony in the region lasted for several centuries.}
{Đây là một danh từ dùng để chỉ quyền bá chủ.}

\vocabentry{Heirloom}
{/ˈerluːm/ (n)}
{A valuable object that has belonged to a family for several generations.}
{This antique watch is a family heirloom that has been passed down for over a hundred years.}
{Đây là một danh từ dùng để chỉ vật gia bảo.}

\vocabentry{Heraldry}
{/ˈherəldri/ (n)}
{The system by which coats of arms and other armorial bearings are devised, described, and regulated.}
{Heraldry was an important part of medieval European culture and social status.}
{Đây là một danh từ dùng để chỉ thuật huy hiệu học.}

\vocabentry{Hero}
{/ˈhɪroʊ/ (n)}
{A person who is admired or idealized for courage, outstanding achievements, or noble qualities.}
{Every culture has its own legendary heroes who are celebrated in songs and stories.}
{Đây là một danh từ dùng để chỉ người hùng dân tộc.}

\vocabentry{Hierarchy}
{/ˈhaɪərɑːrki/ (n)}
{A system or organization in which people or groups are ranked one above the other according to status or authority.}
{There was a strict social hierarchy in feudal societies, with the king at the top.}
{Đây là một danh từ dùng để chỉ hệ thống cấp bậc xã hội.}

\vocabentry{Historiography}
{/hɪˌstɔːriˈɑːɡrəfi/ (n)}
{The study of historical writing.}
{Historiography examines how our understanding of the past has changed over time.}
{Đây là một danh từ dùng để chỉ sử học thành văn hoặc phương pháp viết sử.}

\vocabentry{Icon}
{/ˈaɪkɑːn/ (n)}
{A person or thing regarded as a representative symbol or as worthy of veneration.}
{The ancient church is famous for its beautiful icons depicting saints and biblical scenes.}
{Đây là một danh từ dùng để chỉ biểu tượng văn hóa hoặc tôn giáo.}

\vocabentry{Identity}
{/aɪˈdentəti/ (n)}
{The fact of being who or what a person or thing is.}
{Learning about your national history is essential for developing a strong sense of identity.}
{Trong lịch sử là bản sắc dân tộc/văn hóa.}

\vocabentry{Immigration}
{/ˌɪmɪˈɡreɪʃn/ (n)}
{The action of coming to live permanently in a foreign country.}
{Waves of immigration in the 20th century have made the country more diverse and multicultural.}
{Đây là một danh từ dùng để chỉ sự nhập cư, yếu tố thay đổi lịch sử.}

\vocabentry{Industrialization}
{/ɪnˌdʌstriələˈzeɪʃn/ (n)}
{The development of industries in a country or region on a wide scale.}
{Industrialization in the 19th century led to rapid urban growth and many social problems.}
{Đây là một danh từ dùng để chỉ sự công nghiệp hóa.}

\vocabentry{Inequality}
{/ˌɪnɪˈkwɑːləti/ (n)}
{Difference in size, degree, circumstances, etc.; lack of equality.}
{Historians study the long-term effects of social and economic inequality on societies.}
{Bất bình đẳng xã hội trong lịch sử.}

\vocabentry{Influence}
{/ˈɪnfluəns/ (n/v)}
{The capacity to have an effect on the character, development, or behavior of someone or something.}
{The influence of Roman law can still be seen in many modern legal systems.}
{Đây là một danh từ/động từ chỉ sự ảnh hưởng.}

\vocabentry{Inhabitant}
{/ɪnˈhæbɪtənt/ (n)}
{A person or animal that lives in or occupies a place.}
{The island's original inhabitants were forced to flee when the invaders arrived.}
{Đây là một danh từ dùng để chỉ cư dân hoặc người cư ngụ.}

\vocabentry{Inquiry}
{/ɪnˈkwaɪəri/ (n)}
{An act of asking for information.}
{The historian launched an inquiry into the causes of the mysterious fire.}
{Trong sử học là cuộc điều tra hoặc nghiên cứu lịch sử.}

\vocabentry{Insignia}
{/ɪnˈsɪɡniə/ (n)}
{A badge or distinguishing mark of military rank, office, or membership of an organization.}
{The soldier's uniform bore the insignia of his regiment.}
{Đây là một danh từ dùng để chỉ phù hiệu hoặc cấp hiệu.}

\vocabentry{Institutional}
{/ˌɪnstɪˈtuːʃənl/ (adj)}
{Relating to an established law, practice, or custom.}
{Many historical problems were caused by systemic institutional corruption.}
{Đây là một tính từ dùng để chỉ tính định chế/thể chế.}

\vocabentry{Insurrection}
{/ˌɪnsəˈrekʃn/ (n)}
{A violent uprising against an authority or government.}
{The army was called in to put down the insurrection in the southern province.}
{Đây là một danh từ dùng để chỉ cuộc bạo động hoặc khởi nghĩa vũ trang.}

\vocabentry{Intangible heritage}
{/ɪnˈtændʒəbl ˈherɪtɪdʒ/ (n)}
{Traditions or living expressions inherited from our ancestors and passed on to our descendants.}
{Traditional music and dance are important forms of intangible heritage.}
{Đây là một cụm danh từ dùng để chỉ di sản văn hóa phi vật thể.}

\vocabentry{Integration}
{/ˌɪntɪˈɡreɪʃn/ (n)}
{The action or process of integrating.}
{The integration of new immigrants into the local community is a slow and complex process.}
{Sự hội nhập văn hóa/lịch sử.}

\vocabentry{Interdependence}
{/ˌɪntərdɪˈpendəns/ (n)}
{The dependence of two or more people or things on each other.}
{Modern history is marked by the increasing economic and political interdependence of nations.}
{Sự phụ thuộc lẫn nhau giữa các quốc gia trong lịch sử.}

\vocabentry{Interpretation}
{/ɪnˌtɜːrprɪˈteɪʃn/ (n)}
{The action of explaining the meaning of something.}
{Different historians can have very different interpretations of the same historical event.}
{Sự diễn giải lịch sử.}

\vocabentry{Invasion}
{/ɪnˈveɪʒn/ (n)}
{An instance of invading a country or region with an armed force.}
{The invasion of Poland in 1939 marked the beginning of World War II.}
{Đây là một danh từ dùng để chỉ cuộc xâm lược.}

\vocabentry{Journal}
{/ˈdʒɜːrnl/ (n)}
{A newspaper or magazine that deals with a particular subject.}
{The soldier's personal journal provides a moving account of life in the trenches.}
{Trong lịch sử là nhật ký hoặc tạp chí sử học.}

\vocabentry{Justice}
{/ˈdʒʌstɪs/ (n)}
{Just behavior or treatment.}
{Many people are still seeking justice for the victims of the former regime.}
{Công lý lịch sử.}

\vocabentry{Kinship}
{/ˈkɪnʃɪp/ (n)}
{Blood relationship.}
{The study of kinship patterns is an important part of historical anthropology.}
{Đây là một danh từ dùng để chỉ quan hệ họ hàng hoặc sự gần gũi về nguồn gốc.}

\vocabentry{Knight}
{/naɪt/ (n)}
{(In the Middle Ages) a man who served his sovereign or lord as a mounted soldier in armor.}
{Knights were expected to follow a strict code of conduct called chivalry.}
{Đây là một danh từ dùng để chỉ hiệp sĩ.}

\vocabentry{Knowledge}
{/ˈnɑːlɪdʒ/ (n)}
{Facts, information, and skills acquired by a person through experience or education.}
{A deep knowledge of the past is essential for understanding the present.}
{Tri thức lịch sử.}

\vocabentry{Landmark}
{/ˈlændmɑːrk/ (n)}
{An object or feature of a landscape or town that is easily seen and recognized from a distance.}
{The ancient temple is a famous landmark that attracts thousands of tourists every year.}
{Đây là một danh từ dùng để chỉ địa danh hoặc mốc lịch sử quan trọng.}

\vocabentry{Legend}
{/ˈledʒənd/ (n)}
{A traditional story sometimes popularly regarded as historical but unauthenticated.}
{The legend of King Arthur has been told and retold for centuries.}
{Đây là một danh từ dùng để chỉ truyền thuyết.}

\vocabentry{Legitimacy}
{/ləˈdʒɪtɪməsi/ (n)}
{Conformity to the law or to rules.}
{The new king struggled to establish his legitimacy after the controversial election.}
{Trong chính trị là tính chính danh/hợp pháp của một vương triều.}

\vocabentry{Liberalism}
{/ˈlɪbrəlɪzəm/ (n)}
{A political and social philosophy that promotes individual rights, civil liberties, and democracy.}
{Liberalism was one of the dominant political ideologies of the 19th century.}
{Đây là một danh từ dùng để chỉ chủ nghĩa tự do.}

\vocabentry{Lineage}
{/ˈlɪnieɪdʒ/ (n)}
{Direct descent from an ancestor; ancestry or pedigree.}
{He can trace his lineage back to a famous general who fought in the revolution.}
{Đây là một danh từ dùng để chỉ dòng dõi hoặc huyết thống.}

\vocabentry{Linguistic}
{/lɪŋˈɡwɪstɪk/ (adj)}
{Relating to language or linguistics.}
{The country's rich linguistic heritage is reflected in its many different dialects.}
{Tính từ thuộc về ngôn ngữ - một phần của di sản.}

\vocabentry{Literacy}
{/ˈlɪtərəsi/ (n)}
{The ability to read and write.}
{The rise of mass literacy in the 19th century led to major social and political changes.}
{Sự biết chữ - mốc phát triển văn minh.}

\vocabentry{Litigation}
{/ˌlɪtɪˈɡeɪʃn/ (n)}
{The process of taking legal action.}
{Historical litigation over land rights can last for many decades.}
{Trong sử học thường chỉ các vụ kiện tụng về đất đai/quyền lợi cổ.}

\vocabentry{Locality}
{/loʊˈæləti/ (n)}
{An area or neighborhood.}
{Each locality has its own unique historical character and traditions.}
{Đây là một danh từ dùng để chỉ địa phương hoặc khu vực cư trú.}

\vocabentry{Logic}
{/ˈlɑːdʒɪk/ (n)}
{Reasoning conducted or assessed according to strict principles of validity.}
{There is a certain historical logic to the way the events unfolded.}
{Logic lịch sử/Lý luận.}

\vocabentry{Longevity}
{/lɔːnˈdʒevəti/ (n)}
{Long life.}
{The longevity of the dynasty was due to its strong central government and efficient bureaucracy.}
{Sự trường tồn của một đế chế hoặc truyền thống.}

\vocabentry{Looting}
{/ˈluːtɪŋ/ (n/v)}
{Steal goods from (a place), typically during a war or riot.}
{Widespread looting of ancient sites has led to the loss of many valuable artifacts.}
{Hành vi cướp phá di sản.}

\vocabentry{Mainstream}
{/ˈmeɪnstriːm/ (adj)}
{Representing widespread current thought.}
{Most mainstream historians agree that the war was caused by economic factors.}
{Lịch sử chính thống.}

\vocabentry{Management}
{/ˈmænɪdʒmənt/ (n)}
{The process of dealing with or controlling things.}
{Effective management of national parks is essential for preserving wildlife and heritage.}
{Quản lý di sản - Heritage Management.}

\vocabentry{Manifesto}
{/ˌmænɪˈfestoʊ/ (n)}
{A public declaration of policy and aims.}
{The party's election manifesto included a promise to increase funding for the arts.}
{Đây là một danh từ dùng để chỉ bản tuyên ngôn chính trị.}

\vocabentry{Manuscript}
{/ˈmænjuskrɪpt/ (n)}
{A book, document, or piece of music written by hand rather than typed or printed.}
{The library has a rare collection of ancient medieval manuscripts.}
{Đây là một danh từ dùng để chỉ bản thảo viết tay cổ.}

\vocabentry{Market}
{/ˈmɑːrkɪt/ (n)}
{An area or arena in which commercial dealings are conducted.}
{The town's historic market has been a center of social and economic life for centuries.}
{Chợ truyền thống/Thương trường lịch sử.}

\vocabentry{Masquerade}
{/ˌmæskəˈreɪd/ (n/v)}
{A false show or pretense; a masked ball.}
{The masquerade was a popular social event among the aristocracy in the 18th century.}
{Lễ hội hóa trang truyền thống.}

\vocabentry{Material}
{/məˈtɪriəl/ (n/adj)}
{The matter from which a thing is or can be made; also physical things.}
{Historians study the material culture of the past to understand how people lived.}
{Văn hóa vật thể/Tư liệu.}

\vocabentry{Measurement}
{/ˈmeʒərmənt/ (n)}
{The action of measuring something.}
{Accurate measurement of ancient ruins is essential for archaeological research.}
{Đo lường thời gian/kích thước công trình cổ.}

\vocabentry{Mechanism}
{/ˈmekənɪzəm/ (n)}
{A system of parts working together in a machine; or a process by which something takes place.}
{The mechanism of power in the ancient empire was based on a complex system of patronage.}
{Cơ chế lịch sử.}

\vocabentry{Mediation}
{/ˌmiːdiˈeɪʃn/ (n)}
{Intervention in a dispute in order to resolve it.}
{The conflict was eventually settled through the mediation of a neutral neighboring state.}
{Sự hòa giải giữa các quốc gia trong quá khứ.}

\vocabentry{Memoir}
{/ˈmemwɑːr/ (n)}
{A historical account or biography written from personal knowledge or special sources.}
{The former President wrote a fascinating memoir about his time in office.}
{Đây là một danh từ dùng để chỉ hồi ký.}

\vocabentry{Memory}
{/ˈmeməri/ (n)}
{The faculty by which the mind stores and remembers information.}
{Collective memory plays a crucial role in shaping a nation's sense of identity.}
{Ký ức lịch sử.}

\vocabentry{Mercantilism}
{/ˈmɜːrkəntɪlɪzəm/ (n)}
{An economic theory that trade generates wealth and is stimulated by the accumulation of profitable balances.}
{Mercantilism was the dominant economic policy in Europe during the 17th century.}
{Chủ nghĩa trọng thương.}

\vocabentry{Metaphor}
{/ˈmetəfɔːr/ (n)}
{A figure of speech in which a word or phrase is applied to an object or action.}
{The fall of the wall served as a powerful metaphor for the end of the Cold War.}
{Phép ẩn dụ trong các câu chuyện lịch sử.}

\vocabentry{Meteorology}
{/ˌmiːtiəˈrɑːlədʒi/ (n)}
{The branch of science concerned with the processes and phenomena of the atmosphere.}
{Historical meteorology helps us understand how past weather patterns affected agriculture.}
{Khí tượng học lịch sử - nghiên cứu khí hậu xưa.}

\vocabentry{Metropolis}
{/məˈtrɑːpəlɪs/ (n)}
{The capital or chief city of a country or region.}
{Ancient Rome was a vast metropolis with a population of over a million people.}
{Cố đô hoặc đô thị cổ vĩ đại.}

\vocabentry{Microcosm}
{/ˈmaɪkrəkɑːzəm/ (n)}
{A community, place, or situation regarded as encapsulating in miniature the characteristic qualities of something much larger.}
{The small village was a microcosm of the social tensions that were affecting the entire country.}
{Hình ảnh thu nhỏ của xã hội lịch sử.}

\vocabentry{Migration}
{/maɪˈɡreɪʃn/ (n)}
{Movement of people to a new area or country.}
{The Great Migration of the 1930s saw millions of people moving from rural to urban areas.}
{Sự di cư trong lịch sử.}

\vocabentry{Militancy}
{/ˈmɪlɪtənsi/ (n)}
{The use of confrontational or violent methods in support of a political or social cause.}
{The rise of political militancy led to an increase in social unrest during the decade.}
{Tính chiến đấu hoặc tinh thần quân phiệt.}

\vocabentry{Minority}
{/maɪˈnɔːrəti/ (n)}
{The smaller number or part.}
{Historians often study the experiences of ethnic and religious minorities in past societies.}
{Sắc tộc/Nhóm thiểu số trong lịch sử.}

\vocabentry{Misconduct}
{/ˌmɪsˈkɑːndʌkt/ (n/v)}
{Unacceptable or improper behavior.}
{The official was dismissed for professional misconduct and corruption.}
{Sai phạm của quan chức lịch sử.}

\vocabentry{Missionary}
{/ˈmɪʃəneri/ (n)}
{A person sent on a religious mission, especially one sent to promote Christianity in a foreign country.}
{European missionaries played a major role in the cultural exchange during the colonial era.}
{Nhà truyền giáo.}

\vocabentry{Mobilization}
{/ˌmoʊbələˈzeɪʃn/ (n)}
{The action of a country or its government preparing and organizing troops for active service.}
{The mobilization of the national guard was a clear sign that war was imminent.}
{Sự huy động quân đội.}

\vocabentry{Modernization}
{/ˌmɑːrdənəˈzeɪʃn/ (n)}
{The process of adapting something to modern needs or habits.}
{The government's modernization program aimed to transform the national economy.}
{Đây là một danh từ dùng để chỉ sự hiện đại hóa.}

\vocabentry{Monastery}
{/ˈmɑːnəsteri/ (n)}
{A building or buildings occupied by a community of monks living under religious vows.}
{The medieval monastery was a center of learning and preserved many ancient texts.}
{Tu viện - trung tâm văn hóa lịch sử.}

\vocabentry{Monolith}
{/ˈmɑːnəlɪθ/ (n)}
{A large single upright block of stone, especially one shaped into or serving as a pillar or monument.}
{The ancient temple features several massive stone monoliths carved with mysterious symbols.}
{Khối đá đơn độc cổ.}

\vocabentry{Monumental}
{/ˌmɑːnjuˈmentl/ (adj)}
{Great in importance, extent, or size.}
{The historian's work was a monumental achievement that changed our view of the past.}
{Có tính chất vĩ đại/tượng đài.}

\vocabentry{Morality}
{/məˈræləti/ (n)}
{Principles concerning the distinction between right and wrong or good and bad behavior.}
{The prevailing morality of the Victorian era was characterized by strict social codes.}
{Đạo đức xã hội thời xưa.}

\vocabentry{Mosaic}
{/moʊˈzeɪɪk/ (n)}
{A picture or pattern produced by arranging together small colored pieces of hard material, such as stone, tile, or glass.}
{The floors of the Roman villa were decorated with beautiful and detailed mosaics.}
{Tranh khảm cổ.}

\vocabentry{Motive}
{/ˈmoʊtɪv/ (n)}
{A reason for doing something, especially one that is hidden or not obvious.}
{Historians debate the true motives behind the decision to declare war.}
{Động cơ của các hành động lịch sử.}

\vocabentry{Movement}
{/ˈmuːvmənt/ (n)}
{A group of people working together to advance their shared political, social, or artistic ideas.}
{The civil rights movement fought for racial equality and justice in the United States.}
{Phong trào lịch sử.}

\vocabentry{Multi-ethnic}
{/ˌmʌltiˈeθnɪk/ (adj)}
{Relating to or involving several ethnic groups.}
{The empire was a vast multi-ethnic state with dozens of different languages and cultures.}
{Tính đa sắc tộc.}

\vocabentry{Municipality}
{/mjuːˌnɪsɪˈpæləti/ (n)}
{A city or town that has corporate status and local government.}
{The municipality was granted a royal charter in the 14th century.}
{Đô thị tự trị lịch sử.}

\vocabentry{Mural}
{/ˈmjʊrəl/ (n)}
{A painting or other work of art executed directly on a wall.}
{Archaeologists discovered ancient murals depicting daily life in a Mayan city.}
{Tranh tường cổ.}

\vocabentry{Museum}
{/mjuˈziːəm/ (n)}
{A building in which objects of historical, scientific, artistic, or cultural interest are stored and exhibited.}
{The British Museum has one of the world's largest collections of ancient artifacts.}
{Bảo tàng.}

\vocabentry{Myth}
{/mɪθ/ (n)}
{A traditional story, especially one concerning the early history of a people.}
{Ancient Greek myths often feature powerful gods and heroic figures.}
{Huyền thoại hoặc thần thoại.}

\vocabentry{Mythology}
{/mɪˈθɑːlədʒi/ (n)}
{A collection of myths, especially one belonging to a particular religious or cultural tradition.}
{Nordic mythology is full of stories about brave warriors and magical creatures.}
{Hệ thần thoại.}

\vocabentry{Narrative}
{/ˈnærətɪv/ (n)}
{A spoken or written account of connected events; a story.}
{The author uses a compelling narrative to bring the historical characters to life.}
{Lời kể/Văn tự sự lịch sử.}

\vocabentry{Nationality}
{/ˌnæʃəˈnæləti/ (n)}
{The status of belonging to a particular nation.}
{He holds dual nationality, as his mother is French and his father is American.}
{Quốc tịch/Dân tộc tính.}

\vocabentry{Native}
{/ˈneɪtɪv/ (n/adj)}
{A person born in a specified place or associated with a place by birth.}
{The museum has a collection of art and artifacts from the native people of Australia.}
{Người bản xứ hoặc thuộc về nơi sinh.}

\vocabentry{Natural resources}
{/ˌnætʃrəl ˈriːsɔːrsɪz/ (n)}
{Materials or substances such as minerals, forests, water, and fertile land that occur in nature.}
{The country's abundant natural resources were a major factor in its rapid economic growth.}
{Tài nguyên thiên nhiên ảnh hưởng đến lịch sử.}

\vocabentry{Navigation}
{/ˌnævɪˈɡeɪʃn/ (n)}
{The process or activity of accurately ascertaining one's position and planning and following a route.}
{Advances in navigation made it possible for explorers to cross the vast oceans.}
{Hàng hải lịch sử.}

\vocabentry{Necessity}
{/nəˈsesəti/ (n)}
{The fact of being required or indispensable.}
{The development of farming was a historical necessity for the growth of early civilizations.}
{Sự tất yếu của lịch sử.}

\vocabentry{Negotiation}
{/nəˌɡoʊʃiˈeɪʃn/ (n)}
{Discussion aimed at reaching an agreement.}
{The two nations entered into long negotiations to resolve the border dispute.}
{Đàm phán lịch sử/ngoại giao.}

\vocabentry{Neighborhood}
{/ˈneɪbərhʊd/ (n)}
{A district, especially one forming a community within a town or city.}
{They live in a historic neighborhood with beautiful old houses and tree-lined streets.}
{Khu phố truyền thống.}

\vocabentry{Neutrality}
{/nuːˈtræləti/ (n)}
{The state of not supporting or helping either side in a conflict.}
{Switzerland is famous for its long history of political and military neutrality.}
{Sự trung lập trong chiến tranh.}

\vocabentry{Nomadic}
{/noʊˈmædɪk/ (adj)}
{Living the life of a nomad; wandering.}
{The ancient Mongol tribes were nomadic people who lived in yurts and herded horses.}
{Tính chất du mục.}

\vocabentry{Non-fiction}
{/ˌnɑːn ˈfɪkʃn/ (n)}
{Writing that is based on facts, real events, and real people, such as biography or history.}
{I enjoy reading non-fiction books about ancient Egyptian history.}
{Văn học phi hư cấu/Sách sử.}

\vocabentry{Norm}
{/nɔːrm/ (n)}
{Something that is usual, typical, or standard.}
{Respect for elders is a fundamental social norm in many Asian cultures.}
{Chuẩn mực xã hội trong quá khứ.}

\vocabentry{Notion}
{/ˈnoʊʃn/ (n)}
{A conception of or belief about something.}
{The notion of personal freedom has changed significantly over the last few centuries.}
{Khái niệm hoặc quan niệm lịch sử.}

\vocabentry{Novel}
{/ˈnɑːvl/ (n/adj)}
{A fictitious prose narrative of book length; or new and unusual.}
{The historical novel is set in London during the time of the Great Plague.}
{Tiểu thuyết lịch sử hoặc sự mới lạ.}

\vocabentry{Nuance}
{/ˈnuːɑːns/ (n)}
{A subtle difference in or shade of meaning, expression, or sound.}
{Historians must be careful to capture the subtle nuances of political debate.}
{Sắc thái tinh tế trong phân tích sử học.}

\vocabentry{Nurture}
{/ˈnɜːrtʃər/ (v)}
{Care for and encourage the growth or development of.}
{It is important to nurture a child's natural interest in history and heritage.}
{Nuôi dưỡng truyền thống văn hóa.}

\vocabentry{Objective}
{/əbˈdʒektɪv/ (adj)}
{(Of a person or their judgment) not influenced by personal feelings or opinions.}
{A historian's goal is to provide an objective account of past events.}
{Tính khách quan trong chép sử.}

\vocabentry{Obligation}
{/ˌɑːblɪˈɡeɪʃn/ (n)}
{An act or course of action to which a person is morally or legally bound.}
{We have a moral obligation to protect our cultural heritage for future generations.}
{Nghĩa vụ của thế hệ sau với di sản.}

\vocabentry{Observation}
{/ˌɑːbzərˈveɪʃn/ (n)}
{The action or process of observing something or someone carefully.}
{Careful observation of the site revealed several previously undiscovered artifacts.}
{Sự quan sát/nghiên cứu thực địa khảo cổ.}

\vocabentry{Obsolete}
{/ˌɑːbsəˈliːt/ (adj)}
{No longer produced or used; out of date.}
{Many traditional farming tools have become obsolete due to modernization.}
{Lỗi thời, không còn sử dụng.}

\vocabentry{Official}
{/əˈfɪʃl/ (n/adj)}
{A person holding public office or having official duties.}
{The official records provide a detailed account of the government's activities.}
{Quan chức hoặc mang tính chính thống.}

\vocabentry{Offspring}
{/ˈɔːfsprɪŋ/ (n)}
{A person's child or children.}
{The king was desperate to have male offspring to secure the future of the dynasty.}
{Hậu duệ/Con cháu.}

\vocabentry{Oligarchy}
{/ˈɑːlɪɡɑːrki/ (n)}
{A small group of people having control of a country, organization, or institution.}
{Ancient Sparta was ruled by a powerful military oligarchy.}
{Chính thể tập quyền tay thiểu số trong quá khứ.}

\vocabentry{Opinion}
{/əˈpɪnjən/ (n)}
{A view or judgment formed about something.}
{In my opinion, the most significant historical event of the 20th century was the moon landing.}
{Quan điểm sử học.}

\vocabentry{Oppression}
{/əˈpreʃn/ (n)}
{Prolonged cruel or unjust treatment or control.}
{The struggle for freedom was a fight against decades of systematic political oppression.}
{Sự áp bức lịch sử.}

\vocabentry{Option}
{/ˈɑːpʃn/ (n)}
{A thing that is or may be chosen.}
{The government had several options for dealing with the crisis, but none were perfect.}
{Các lựa chọn mang tính bước ngoặt lịch sử.}

\vocabentry{Opulence}
{/ˈɑːpjələns/ (n)}
{Great wealth or luxuriousness.}
{The opulence of the royal court was a sharp contrast to the poverty of the common people.}
{Sự sang trọng/xa hoa của các triều đại.}

\vocabentry{Oral}
{/ˈɔːrəl/ (adj)}
{Spoken rather than written.}
{Many folktales were preserved through oral tradition for centuries before being written down.}
{Tính chất truyền miệng.}

\vocabentry{Order}
{/ˈɔːrdər/ (n/v)}
{The arrangement or disposition of people or things.}
{The king issued an order for the immediate mobilization of the army.}
{Trật tự xã hội/Sắc lệnh.}

\vocabentry{Organization}
{/ˌɔːrɡənəˈzeɪʃn/ (n)}
{An organized body of people with a particular purpose.}
{The historical society is a non-profit organization dedicated to preserving local heritage.}
{Tổ chức/Thiết chế xã hội xưa.}

\vocabentry{Orientation}
{/ˌɔːriənˈteɪʃn/ (n)}
{The action of orienting someone or something.}
{The temple was built with a specific orientation towards the rising sun.}
{Sự định hướng văn hóa/tôn giáo.}

\vocabentry{Origin}
{/ˈɔːrɪdʒɪn/ (n)}
{The point or place where something begins, arises, or is derived.}
{Historians are still debating the true origins of the ancient civilization.}
{Nguồn gốc/Cội nguồn.}

\vocabentry{Ornate}
{/ɔːrˈneɪt/ (adj)}
{Made in an intricate shape or decorated with complex patterns.}
{The ancient palace features beautifully ornate carvings and mosaics.}
{Được trang trí công phu, hoa mỹ.}

\vocabentry{Outcome}
{/ˈaʊtkʌm/ (n)}
{The way a thing turns out; a consequence.}
{The outcome of the battle changed the course of the war and the history of the region.}
{Kết quả của một sự kiện lịch sử.}

\vocabentry{Outline}
{/ˈaʊtlaɪn/ (n/v)}
{A general description or plan giving the essential features of something.}
{The historian provided a brief outline of the causes of the revolution.}
{Dàn ý/Đề cương lịch sử.}

\vocabentry{Outpost}
{/ˈaʊtpoʊst/ (n)}
{A small military camp or settlement in a remote area.}
{The remote military outpost was abandoned after the end of the war.}
{Trạm tiền tiêu/Thuộc địa xa xôi.}

\vocabentry{Overcome}
{/ˌoʊvərˈkʌm/ (v)}
{Succeed in dealing with a problem or difficulty.}
{The nation was able to overcome the challenges of the economic depression through hard work.}
{Vượt qua nghịch cảnh lịch sử.}

\vocabentry{Overthrow}
{/ˌoʊvərˈθroʊ/ (v)}
{Remove forcibly from power.}
{The revolutionary group aimed to overthrow the dictator and install a democratic government.}
{Lật đổ chính quyền.}

\vocabentry{Ownership}
{/ˈoʊnərʃɪp/ (n)}
{The act, state, or right of possessing something.}
{The local community is fighting for the legal ownership of the historic land.}
{Quyền sở hữu di sản.}

\vocabentry{Pacifism}
{/ˈpæsɪfɪzəm/ (n)}
{The belief that any violence, including war, is unjustifiable.}
{His commitment to pacifism led him to oppose the government's decision to go to war.}
{Chủ nghĩa hòa bình.}

\vocabentry{Paradigm}
{/ˈpærədaɪm/ (n)}
{A typical example or pattern of something; a model.}
{The discovery of DNA created a new paradigm for biological research and history.}
{Mô hình/Chuẩn mực lịch sử.}

\vocabentry{Parliament}
{/ˈpɑːrləmənt/ (n)}
{The highest legislature of a country.}
{Parliament met to debate the new tax laws proposed by the government.}
{Nghị viện/Quốc hội - thiết chế lịch sử lâu đời.}

\vocabentry{Participation}
{/pɑːrˌtɪsɪˈpeɪʃn/ (n)}
{The action of taking part in something.}
{The mass participation of ordinary people was essential for the success of the movement.}
{Sự tham gia của dân chúng vào các sự kiện lịch sử.}

\vocabentry{Passage}
{/ˈpæsɪdʒ/ (n)}
{A journey by sea or air; or a short part of a book.}
{The dangerous passage across the Atlantic took several weeks by ship.}
{Chuyến hải hành lịch sử hoặc đoạn trích.}

\vocabentry{Pastoral}
{/ˈpæstərəl/ (adj)}
{Portraying or evoking country life, typically in a romanticized way.}
{The artist is known for his idyllic pastoral scenes of grazing sheep and rolling hills.}
{Phong cách điền viên/đồng quê trong sử liệu.}

\vocabentry{Patriarchal}
{/ˌpeɪtriˈɑːrkl/ (adj)}
{Relating to or characteristic of a system of society or government controlled by men.}
{Most ancient civilizations were patriarchal societies with very few rights for women.}
{Tính gia trưởng/phụ hệ trong xã hội xưa.}

\vocabentry{Patriotism}
{/ˈpeɪtriətɪzəm/ (n)}
{Devotion to and vigorous support for one's country.}
{Patriotism can be a powerful force for national unity during times of crisis.}
{Lòng yêu nước.}

\vocabentry{Patronage}
{/ˈpætrənɪdʒ/ (n)}
{The support given by a patron to an artist or organization.}
{Without the patronage of the wealthy family, many Renaissance works would not exist.}
{Sự bảo trợ nghệ thuật/lịch sử.}

\vocabentry{Pedagogy}
{/ˈpedəɡɑːdʒi/ (n)}
{The method and practice of teaching.}
{Modern pedagogy has been influenced by many historical educational movements.}
{Nghiệp vụ sư phạm/Lịch sử giáo dục.}

\vocabentry{Penalty}
{/ˈpenəlti/ (n)}
{A punishment imposed for breaking a law, rule, or contract.}
{In ancient times, the penalty for theft was often very severe.}
{Hình phạt trong luật pháp cổ.}

\vocabentry{Periodical}
{/ˌpɪriˈɑːdɪkl/ (n)}
{A magazine or newspaper published at regular intervals.}
{The library has a large collection of 19th-century periodicals and journals.}
{Ấn phẩm định kỳ lịch sử.}

\vocabentry{Perjury}
{/ˈpɜːrdʒəri/ (n)}
{The offense of willfully telling an untruth in a court.}
{Several witnesses were accused of perjury during the high-profile treason trial.}
{Tội khai man trước tòa trong sử liệu.}

\vocabentry{Persistence}
{/pərˈsɪstəns/ (n)}
{Firm or obstinate continuance in a course of action.}
{The persistence of traditional customs is a testament to the strength of the community.}
{Sự kiên trì gìn giữ văn hóa.}

\vocabentry{Perspective}
{/pərˈspektɪv/ (n)}
{A particular attitude toward or way of regarding something.}
{Traveling to different countries can provide a new perspective on your own national history.}
{Góc nhìn lịch sử.}

\vocabentry{Petitioner}
{/pəˈtɪʃənər/ (n)}
{A person who presents a petition to an authority.}
{A group of petitioners gathered at the palace to ask the king for lower taxes.}
{Người kiến nghị/thỉnh nguyện trong lịch sử.}

